\documentclass{article}

\usepackage[margin=1in]{geometry}
\linespread{1.2}
\usepackage{amsmath,amsthm,amssymb,amsfonts,commath}


\newcommand{\mto}{\stackrel{\mu}{\rightarrow}}
\newcommand{\Mcal}{\mathcal{M}}


 
\title{Real Analysis I --- Homework 4}

% Replace with your name and ID below
\author{FirstName LastName \\ ID: 123-45-6789}

\date{}


\begin{document}

\maketitle

Unless otherwise specified, all functions are defined on $E \subset \mathbb{R}^{n}$ where $E \in \mathcal{M}$.
\bigskip

For notation simplicity, a set such as $\{x \in E: f(x) > t\}$ is written as $\{f>t\}$ for short.

\begin{enumerate}

\item Suppose $f_{k} \mto f$ on $E$, and $f_{k} \le g$ a.e.~$E$. Show that $f \le g$ a.e.~$E$.

%uncomment below and write your proof after it.
%\begin{proof}
%\end{proof}


%\newpage % uncomment this line to start on a new page

\item Suppose $f_{k} \mto f$ on $E$, and $f_{k} \le f_{k+1}$ a.e.~$E$ for all $k \in \mathbb{N}$. Show that $f_{k}\to f$ a.e.~$E$.

%uncomment below and write your proof after it.
%\begin{proof}
%\end{proof}


%\newpage % uncomment this line to start on a new page

\item Find the integral $\int_{0}^{1} f$ where 
\[
f(x) = 
\begin{cases}
\frac{1}{\sqrt{x}} & \mbox{if } x \notin \mathbb{Q} \\
x^{3} & \mbox{if } x \in \mathbb{Q} 
\end{cases}
\]

%uncomment below and write your proof after it.
%\begin{proof}
%\end{proof}


%\newpage % uncomment this line to start on a new page

\item Suppose $f, g \in L(E)$. Prove that $\sqrt{|f|^{2}+|g|^{2}} \in L(E)$.

%uncomment below and write your proof after it.
%\begin{proof}
%\end{proof}


%\newpage % uncomment this line to start on a new page

\item Suppose $f\in L(E)$, and for any bounded measurable function $\phi: E \to \mathbb{R}$ there is $\int_{E} f\phi = 0$. Show that $f=0$ a.e.~$E$.

%uncomment below and write your proof after it.
%\begin{proof}
%\end{proof}


%\newpage % uncomment this line to start on a new page

\item Suppose $f\in L(E)$ and $\mu(E) < \infty$. Show that for any $\epsilon > 0$, there exists a bounded measurable function $g: E \to \mathbb{R}$ such that $\int_{E} |f-g| < \epsilon$.

%uncomment below and write your proof after it.
%\begin{proof}
%\end{proof}


%\newpage % uncomment this line to start on a new page


\item Suppose $\mu(E)< \infty$. Show that $f \in L(E)$ iff $\sum_{k=0}^{\infty} \mu(\{|f| \ge k\}) < \infty$.

%uncomment below and write your proof after it.
%\begin{proof}
%\end{proof}



%\newpage % uncomment this line to start on a new page


\item Prove that for any $p > -1$ there is
\[
\int_{0}^{1} \frac{x^{p}}{1-x} \ln \frac{1}{x} \dif x = \sum_{k=1}^{\infty} \frac{1}{(p+k)^{2}}
\]

%uncomment below and write your proof after it.
%\begin{proof}
%\end{proof}


%\newpage % uncomment this line to start on a new page


\item Prove that 
\[
\lim_{k \to \infty} \int_{0}^{\infty} \frac{1}{(1+ \frac{t}{k})^{k} \cdot t^{1/k}} \dif t = 1
\]
%uncomment below and write your proof after it.
%\begin{proof}
%\end{proof}



%\newpage % uncomment this line to start on a new page


\item Suppose $f$ is integrable on $(0,\infty)$. Define 
\[
g(x) = \int_{0}^{\infty} \frac{f(t)}{x+t} \dif t
\] 
for any $x \in (0,\infty)$. Is $g$ continuous on $(0,\infty)$? Justify your answer.

%uncomment below and write your proof after it.
%\begin{proof}
%\end{proof}




\end{enumerate}


\end{document}